% 第 3 章 关键技术方案 (目标 6 页)
% 对应官方要求: 技术路线 + 算法设计
% 结构: 总体路线 → 规则引擎 → LLM评审 → 协议适配 → 审计 → 工作流程
\section{关键技术方案}

\subsection{总体技术路线}

系统采用三层混合防护,自下而上为结构化规则、结构化 LLM 评审、真实环境接入验证。三层独立可部署,联合提供从已知模式拦截到语义兜底再到端到端验证的完整链路。

第一层结构化规则针对已知可枚举的高风险动作模式,覆盖邮箱外发白名单、写入路径白名单、SQL 写操作、定时任务、审计日志清理等。该层在所有调用链路上最先执行,执行开销极低,可解释性最强。

第二层结构化 LLM 评审针对规则未覆盖但语义可疑的工具调用,在注入识别基础上新增过度代理检测与意图对齐判断两个维度,要求评审模型输出结构化分数与简短中文理由。该层在规则未命中时按需触发。

第三层真实环境接入验证针对链路真实性问题,在真实智能体上跑通端到端攻击触发与拦截,证明防护在真实工具链路生效。

设计原则四条:真实优先,全部核心评估在真实业务话术样本上完成;分层独立,任意一层失能不影响其他层;封闭失败,评审与上游异常一律阻断并升级人工审批;可解释优先,规则给编号与触发条件,评审给分数与理由,审计给完整轨迹。

\subsection{结构化规则引擎设计}

规则引擎基于声明式策略匹配实现。每条规则包含编号、标题、工具选择器、谓词、效果、消息、启用标志七个字段。工具选择器限定规则适用的工具集合,谓词基于参数字段做任意嵌套匹配,效果分阻断、审批、记录三类。

实际部署 14 条规则,覆盖六类高风险动作。代表性规则见表~\ref{tab:rules},完整列表见附录 A。

\begin{table}[H]
  \centering
  \caption{代表性结构规则}
  \label{tab:rules}
  \small
  \begin{tabular}{p{2.8cm} p{2.4cm} p{7.2cm}}
    \toprule
    规则编号 & 工具范围 & 触发条件摘要 \\
    \midrule
    GOV-APPROVE-002 & send\_email 类 & 收件人或抄送含外部、群发、跨部门地址关键字 \\
    GOV-DENY-007 & write 类 & 写入路径含正式、机密、登记、归档目录或版本化文件名 \\
    GOV-DENY-008 & exec 类 & 命令含 SQL 写操作或审批字段改写 \\
    GOV-DENY-009 & cron 类 & 任意定时任务调用,覆盖延迟外发攻击 \\
    GOV-DENY-010 & exec 类 & 命令含日志清理、审计清理、服务重启 \\
    GOV-DENY-012 & write 类 & 参数字段名变体下的敏感路径,覆盖字段命名差异 \\
    \bottomrule
  \end{tabular}
\end{table}

谓词引擎支持约 25 种操作符,含字段存在性、子串匹配、正则匹配、集合成员、数值比较、逻辑组合。新增高风险模式仅需新增一条规则声明,无需改动引擎代码。

\subsection{结构化 LLM 评审设计}

评审层输入为当前工具调用名、参数摘要、最近 5 步前置调用文本。评审模型按 0 到 10 评分,6 分以上触发阻断。

评分沿三个维度展开。第一维度为注入风险,涵盖社工冒充、显式注入、编码绕过、函数返回值注入、知识库投毒、跨会话投毒。第二维度为过度代理,涵盖用户未明确请求的写入加外发组合、群发操作、批量批准、汇总后外发。第三维度为意图对齐,从用户原始请求出发评估当前调用是否在用户意图之内。

评审指令内含若干硬性约束。涉及执行命令、发送邮件、写文件等高风险工具且非用户明确请求时,基础分不低于 6。前置文本含社工信号时基础分不低于 7。前置轨迹出现连续多次同类工具调用且非用户批量请求时基础分不低于 7。完整评审指令见附录 B。

工程安全方面,评审调用结果做封闭失败处理。上游 HTTP 错误、网络异常、响应解析失败一律按阻断处理,不放行。该设计直接消除评审失效即全放行的隐患。

参数提取与降级。评审默认读取响应正文字段。当上游模型将推理结果放入推理字段且正文为空时,自动降级到推理字段。该降级保障多模型兼容,不引入新失败路径。

\subsection{响应解析与协议适配}

接入层解决代理能否拦到真实工具调用的问题。智能体默认发起流式请求,模型上游在流式模式下返回分片事件而非单一 JSON,代理无法在响应解析阶段检查工具调用。代理转发时强制将流式字段改为关闭,上游返回标准 JSON,代理在响应解析阶段提取工具调用并执行策略检查。该改造不影响模型输出语义,仅影响代理侧的解析能力。

平台适配涉及两处。其一,部分桌面操作系统对智能体状态目录施加系统级保护,需通过环境变量将状态目录隔离到可写路径。其二,代理需兼容多家模型上游的协议差异,含正文与推理字段的位置差异。

\subsection{审计溯源设计}

审计层对每一次模型调用记录七类字段:调用时间、目标模型、请求摘要、工具调用清单、命中规则编号、评审分数与理由、最终动作。审计日志以 JSON Lines 格式追加写入,单条记录自包含,支持按样本编号、规则编号、时间段做检索与回放。审计样例见附录 D。

\subsection{系统工作流程}

一次完整的攻击拦截流程包含七个环节,全链路在防护代理内完成,见图~\ref{fig:workflow}。

% TODO: 图 3 工作流程图(七环节泳道: 请求接入 → 响应解析 → 文本级阻断 → 规则层 → 评审层 → 回写响应 → 审计)
\begin{figure}[H]
  \centering
  \fbox{\parbox[c][7cm][c]{0.9\textwidth}{\centering\small 待绘制: 七环节工作流程泳道图}}
  \caption{攻击拦截工作流程}
  \label{fig:workflow}
\end{figure}

\begin{enumerate}[nosep]
  \item 请求接入。智能体发起模型调用,请求经代理转发,代理解析请求体并识别目标模型。
  \item 响应解析。代理强制请求为非流式,上游返回标准 JSON,代理提取消息与工具调用,还原完整调用轨迹。
  \item 文本级阻断。对响应文本做注入标记与高危关键词检测,命中时改写响应并记录审计。
  \item 规则层判定。对每个工具调用执行 14 条结构性规则,任意规则命中即阻断,返回结构化阻断消息并写审计。
  \item 评审层兜底。规则未命中的调用进入评审层,综合当前调用、前置轨迹、用户原始意图三维度评分,6 分以上触发阻断,评审异常时按封闭失败处理。
  \item 回写响应。通过检查的调用原样保留,被阻断的调用替换为阻断提示,智能体据此生成合规回复。
  \item 审计溯源。每个环节的判定结果、命中规则、评分、原始调用与改写内容全部落盘。
\end{enumerate}

以伪造签字类样本为例。用户以业务话术诱导智能体批量出具同意意见,智能体调用命令执行工具发起 SQL 更新。代理在第三环节捕获命令文本,第四环节命中 SQL 写操作规则,响应被改写为阻断提示,审计日志记录完整调用链。全程无需智能体侧任何改造。
